\subsection{Integral flat twists and the passage to an arbitrary weight}
\label{subsec:ftb-flat-twists-weight-reduction}

A componentwise twisting section in this argument consists of
opposite-weight regular rules \(t,t^{-1}\), supported on one
component \(Z\), with \(tt^{-1}=e_Z\).
Multiplication by either rule preserves the indicated cusp spaces.
Its normalized pointwise norm on \(Z\) is a positive constant
\(c_\sigma\) at each coefficient-field embedding.
Its valuations satisfy
\(\rho_w\le v(t)\le\rho'_w\) at every valued point over \(w\);
both functions have finite support.
These are all the fields of the source's \srcid{TwSec}.
In particular cusp preservation and the equality of pointwise
norms are part of this data; they are not inferred solely from the
existence of a rational section.
Its lower degree is
\[
 \deg_{\rm lo}(t)=
   -\frac1{[K:\mathbb Q]}\sum_{\sigma:K\hookrightarrow\mathbb C}
                           \log c_\sigma
   +\sum_w\rho_w\,\frac{\log Nw}{[K:\mathbb Q]} .
\]
The product formula makes this unchanged by rescaling \(t\) by
\(a\in K^\times\), with the corresponding changes in \(c,\rho\).

\begin{srlemma}[Rounding a flat line by one algebraic scalar]
\label{lem:ftb-round-flat-line}
Let \(K\) be a number field of degree \(d_K\).
Let \(c_\sigma>0\) be conjugation invariant, and let \(\rho_w\)
be real, finitely supported and integral at primes not above a
finite set \(\mathcal B\) of rational primes.
Then there is \(a\in K^\times\) with
\[
 v_w(a)+\rho_w\ge0\quad\hbox{for every }w
\]
and, for every \(\sigma\),
\[
 \log(c_\sigma|\sigma a|)
 \le-\deg(c,\rho)+\tfrac12\log\operatorname{rd}K
                      +\sum_{p\in\mathcal B}\log p+1 .
\]
Here \(\deg(c,\rho)\) is the degree in the preceding display.
\end{srlemma}
\begin{proof}
Set \(b=\prod_{p\in\mathcal B}p\), and let \(S\) contain the support
of \(\rho\).
Use the nonzero fractional ideal
\[
                 J=(b)\prod_{w\in S}w^{\lfloor-\rho_w\rfloor}.
\]
If \(a\ne0\) belongs to \(J\), then
\(v_w(a)\ge v_w(b)+\lfloor-\rho_w\rfloor\).
Off \(\mathcal B\), \(\rho_w\) is integral, so this is at least
\(-\rho_w\).
At a prime over \(\mathcal B\), \(v_w(b)\ge1\);
\(1+\lfloor-\rho_w\rfloor>-\rho_w\) proves the same conclusion.
Moreover
\[
 \log NJ
  =d_K\log b+\sum_{w\in S}\lfloor-\rho_w\rfloor\log Nw
  \le d_K\left(\log b-
               \sum_w\rho_w\frac{\log Nw}{d_K}\right).
\]

Apply Minkowski to the real intervals and complex discs
\(|x_\sigma|<R/c_\sigma\), where
\[
 R=\frac52
       \left(NJ\,\sqrt{|\operatorname{disc}K|}
                   \prod_\sigma c_\sigma\right)^{1/d_K}.
\]
This is a symmetric convex body in the mixed embedding.
The covolume of \(J\) is
\(2^{-r_2}NJ\sqrt{|\operatorname{disc}K|}\).
The body's interval/disc factor is at least one in the source's
volume normalization, whereas \((5/2)^{d_K}>2^{d_K}\);
hence its volume is strictly larger than \(2^{d_K}\) times that
covolume.
It contains a nonzero \(a\in J\).
Thus \(c_\sigma|\sigma a|<R\) simultaneously.
Taking logarithms, using \(\log(5/2)<1\), and inserting the ideal
norm bound gives the asserted inequality.
This proof retains the source's fractional-ideal floor and the
extra factor \(\prod_{p\in\mathcal B}p\); it does not round at
different embeddings separately.
\end{proof}

For a function \(w:\Psi\to\mathbb Z\), let
\(\operatorname{par}(w)\) denote the weight constant with value
\(w_\tau\) on every row at the CM embedding \(\tau\).
The actual flat-section packet used here supplies a twisting
section of weight
\(\operatorname{wt}(\operatorname{par}(w),0)\), with integral
valuations off \(\mathcal B\), and lower degree
\begin{equation}\label{eq:ftb-flat-twist-exact-degree}
 \deg_{\rm lo}(t)=
 \sum_{\tau\in\Psi}w_\tau c_1
 -q_{\rm sig}(\log\operatorname{rd}E-\log d_{\rm host})
                                      \sum_{\tau\in\Psi}w_\tau .
\end{equation}
This is the source construction
\srcid{ExistsTwSecOfConst.twSec\_of\_flat}: tensor the actual
rank-one constant sections with exponents \(w_\tau\), and add
the finite defect
\(\sum_\tau(-w_\tau)\operatorname{flDefect}_\tau\).
The normalized sum of each defect is
\(q_{\rm sig}(\log\operatorname{rd}E-\log d_{\rm host})\);
finite support permits exchanging the sums.
The source's host discriminant bound is
\[
 |\log\operatorname{rd}E-\log d_{\rm host}|
                                      \le\log\operatorname{rd}E .
\]
It follows from \eqref{eq:ftb-flat-twist-exact-degree} that
\[
 -\deg_{\rm lo}(t)\le
 (|c_1|+q_{\rm sig}\log\operatorname{rd}E)\sum_\tau|w_\tau|.
\]
The flat packet and its integral defect are the exact geometric
input here.
The condition \srcid{UnramPts} ensures that the valuations of
these constant rules are integral away from \(\mathcal B\), not
merely rational: it supplies an integral valued point of
ramification index one off that set.
Applying Lemma~\ref{lem:ftb-round-flat-line} and multiplying by its
scalar gives an integral twisting section with
\begin{equation}\label{eq:ftb-flat-twist-short}
 \log c_\sigma(t)\le
 (|c_1|+q_{\rm sig}\log\operatorname{rd}E)\sum_\tau|w_\tau|
 +\sum_{p\in\mathcal B}\log p+\tfrac12\log\operatorname{rd}K+1 .
\end{equation}
The degree \(d_K\) cancels before this bound; no polynomial
estimate in a varying auxiliary degree is substituted for it.

The other twist is the fourth-power theta key unit.
Its exact source input is an invertible rule \(K_4\) of weight
\(\operatorname{wt}(0,-4)\), together with an actual key section,
such that
\begin{align}
 0\le v(K_4)&\le
       4\sum_\tau r_\tau
           \operatorname{ord}_w(\tau\mathfrak D_E\mathcal O_K),
                                                   \label{eq:ftb-key-unit-finite}\\
 \operatorname{homN}_\sigma(K_4)&=(2\pi)^{2fr}.
                                                   \label{eq:ftb-key-unit-infinite}
\end{align}
These are exactly the weight-four scalar and two-frame
comparisons of Section~\ref{sec:key-different}.
For the analytic value, at an orthonormal frame the key formula
gives
\[
 |\kappa|=(2\pi)^{2g}(2\pi\xi)^{2g},\qquad g=fr .
\]
Unimodularity of the polarization says that the fundamental
volume times \(\det H\) is one; the theta generator's normalized
square integral is one.
The weight factor is \((2\pi\xi)^{-2g}\), so it cancels and gives
\eqref{eq:ftb-key-unit-infinite}.
Moving a general frame to an orthonormal one preserves its
normalized norm by the weight transformation law.
At a finite place, if \(D\) is the determinant of the complete
integral eigenframe in an integral Lie basis, then
\(K_4=D^4u\), \(u\in V^\times\).
Trace-dual projectors give
\(D\mid\prod_\tau\tau(\delta_\tau)^{r_\tau}\) for elements
\(\delta_\tau\) attaining the local different valuation.
This gives \eqref{eq:ftb-key-unit-finite}.

Consequently \(t_1=e_ZK_4^m\), with inverse \(e_ZK_4^{-m}\),
is a twisting section of weight \(\operatorname{wt}(0,-4m)\),
with nonnegative lower valuations and
\[
                \log c_\sigma(t_1)\le m\,2fr\log(2\pi).
\]
The inverse and all powers are on the same component; their
product is \(e_Z\), not \(1\) on the other components.

\begin{srlemma}[Multiplication transports an integral spanning family]
\label{lem:ftb-twist-span}
Let \(t\) be a twisting section with nonnegative lower valuations.
If \(\mathcal S\) is an integral \(K\)-spanning set of cusp forms
of even weight \(G\), then the products \(e_Z f t\),
\(f\in\mathcal S\), span the component cusp space of weight
\(G+\operatorname{wt}(t)\), are integral, and satisfy
\[
       \|e_Zft\|_{\sigma,Z}
                          \le c_\sigma(t)\,\|f\|_{\sigma,Z}.
\]
\end{srlemma}
\begin{proof}
Multiplication by \(t\) has inverse multiplication by its specified
inverse, because cusp preservation holds both ways and
\(e_Z\) acts as the identity on this component.
It is therefore a \(K\)-linear isomorphism.
The component projection \(f\mapsto e_Zf\) is surjective.
Their composite transports the spanning set.

At every valuation over \(w\), the product has valuation at least
\(v(f)+\rho_w(t)\ge0\).
The integral lattice is defined by these same inequalities.
For the norm, on the component the pointwise norm multiplies
exactly by the positive constant \(c_\sigma(t)\); off the
component the product is zero.
Integrate its square with the source's component measure and
take square roots.
This gives the displayed inequality (in fact equality whenever
both sides use precisely the component projection).
Finiteness is legitimate because a cusp form's pointwise norm is
bounded and the component measure has mass at most one.
The norm and integral lattice used after multiplication are
therefore those of the target weight on the same model.
\end{proof}

\begin{srproposition}[The source weight reduction and short span]
\label{prop:ftb-arbitrary-weight-short-span}
Let the good-cover packet be the one in
Proposition~\ref{prop:ftb-integral-span-assembly}.
Put \(X=frL\), \(f,r,L\ge1\), and assume
\(\log(fr)\le L\).  Retain positive signature
\(p_{\rm sig},q_{\rm sig}\ge1\), \(r=p_{\rm sig}+q_{\rm sig}\),
from that good-cover packet.
Suppose its field-size bound is \(o\le391X^2\), and suppose
\[
 \log\operatorname{rd}E\le3L,\quad
 |c_1|\le9rL,\quad
 \sum_{p\in\mathcal B}\log p\le C_BX^2 .
\]
Retain the actual flat/key twisting sections just described.
Let \(\lambda_{\tau,1}\ge\cdots\ge\lambda_{\tau,r}\) have
signature gap at least \(8fr+10\), and put
\(S=\sum_{\tau,i}|\lambda_{\tau,i}|\), \(\log S\le10L\).
If the even-domain integral spanning bound has constant \(C_G\),
then the component cusp lattice of weight
\(\operatorname{wt}(\lambda,0)\) has an integral \(K\)-spanning
family with every norm at most
\[
                  \exp\{C_2(X^3S+X^5)\},
                   \qquad C_2=11C_G+C_B+500 .
\]
The constants are chosen before all host, component and weight
data.
\end{srproposition}
\begin{proof}
Set \(m=fr+1\) and \(n=8m\).
For each CM embedding set
\[
       n_{0,\tau}=1+4m-\lambda_{\tau,p_{\rm sig}},\qquad
       \lambda'=\lambda+\operatorname{par}(n_0).
\]
Every entry in the first signature block is strictly greater
than \(4m\).
The gap condition says the first entry of the second block is
at most \(-4m-1\).
Antitonicity gives this bound for the whole second block.
Thus the following nonnegative integer rows define theta
exponents:
\[
 a_{\tau,i}=\lambda'_{\tau,i}-4m
                    \quad(1\le i\le p_{\rm sig}),\qquad
 a_{\bar\tau,j}=-\lambda'_{\tau,r+1-j}-4m
                    \quad(1\le j\le q_{\rm sig}),
\]
with zero continuation outside these finite rows.
They are antitone, and \(n\) is even, so
\(\operatorname{Dom}(a,n)\) holds literally.
The first-block exponents are positive, as required for the
cuspidal theta spanning theorem.
The associated weight is \(\operatorname{laOf}(a,n)=\lambda'\).
Also \(n=8(fr+1)\ge8fr+4\), the required column threshold.

The triangle inequality gives
\[
 N_0:=\sum_\tau|n_{0,\tau}|
                    \le S+f(4(fr+1)+1),\qquad
 \kappa:=\operatorname{sizeK}(a,n)
                    \le(1+r)S+rf(4(fr+1)+1).
\]
The second follows from the identity
\(\kappa=\sum_{\tau,i}|\lambda_{\tau,i}+n_{0,\tau}|\).
These estimates and \(\log S\le10L\), \(\log(fr)\le L\)
also give \(\log\kappa\le20L\).
For instance \(S\ge1\) and \(fr\ge1\), so
\(\kappa\le 11(fr)^2(S+1)\);
we have
\(\log\kappa\le\log22+2\log(fr)+\log S
               \le\log22+12L\le20L\).
Here \(S\ge1\), and \(\log22<8\).
The same estimate is the source's integral-weight logarithm
calculation.
Apply the even-domain result: there is an integral spanning
family at weight
\(\operatorname{wt}(\lambda',4m)\), of norm
\(\le\exp(C_GX^2\kappa)\).
The identification of its section and Petersson norm with
the theta-exponent notation uses the exact identity of the full
exponent function; it retains the evenness of \(n\).

Multiply by \(t_1=e_ZK_4^m\) and by the integral flat twist \(t_2\)
of parallel weight \(-n_0\).
Their weights add exactly to
\[
 \operatorname{wt}(\lambda',4m)+
 \operatorname{wt}(0,-4m)+
 \operatorname{wt}(\operatorname{par}(-n_0),0)
                                  =\operatorname{wt}(\lambda,0).
\]
Lemma~\ref{lem:ftb-twist-span} proves both spanning and integrality.
The logarithm of every resulting norm is at most
\begin{align*}
 &(fr+1)2fr\log(2\pi)\\
 &\quad+(|c_1|+q_{\rm sig}\log\operatorname{rd}E)N_0
       +C_BX^2+\tfrac12\log\operatorname{rd}K+1
       +C_GX^2\kappa .
\end{align*}
Here \(\log\operatorname{rd}K\le o\le391X^2\).
We give the actual scalar absorption.
Since \(fr\le X\), \(rL\le X\), \(X\ge1\), and
\(\log(2\pi)\le7\), the first term is at most \(28X^5\).
The middle product is at most
\[
                  12X^3S+108X^5 ,
\]
because \(|c_1|+q_{\rm sig}\log\operatorname{rd}E\le12rL\)
and \(N_0\le S+9f(fr)\).
Finally
\[
    \kappa\le2XS+9X^2,\qquad
    C_GX^2\kappa\le2C_GX^3S+9C_GX^5.
\]
The field term is at most \(196X^5\), \(1\le X^5\), and
\(C_BX^2\le C_BX^5\).
Adding proves the claimed bound with \(11C_G+C_B+500\).
This is the source's
\srcid{CovolumeNumericalBoundsCollect.covol\_numerical\_collect}.
\end{proof}

\begin{srproposition}[Short generators, with the same window variables]
\label{prop:ftb-arbitrary-weight-short-generators}
In the preceding setting, suppose the component cusp lattice
\(L_\lambda\) is a full \(\mathcal O_K\)-lattice, its Petersson
norms arise from the same positive Gram frame, and its
\(K\)-rank \(h\) satisfies \(\log h\le C_bX^3\).
Then there is a family of \(\mathbb Z\)-generators of \(L_\lambda\)
with common norm bound \(A\ge1\), and cardinality \(n_{\rm gen}\),
such that
\[
 \log A\le C(X^3S+X^5),\qquad
 \log n_{\rm gen}\le CX^3,\qquad
                       C=2000+C_b+C_2 .
\]
The zero-rank lattice has the empty generating family, with
\(A=1\); the second inequality is then interpreted by this
separate convention.
\end{srproposition}
\begin{proof}
Apply Proposition~\ref{prop:ftb-short-generators} to the integral
spanning family:
\[
 n_{\rm gen}\le2d_Kh,\qquad
 A\le\max\{1,d_K^4h\operatorname{rd}K
                         \exp(C_2(X^3S+X^5))\}.
\]
The norm hypotheses in that proposition hold for these exact
Petersson norms: scaling has factor \(|\sigma a|\), nonnegativity
is immediate, and the triangle inequality is the Euclidean
triangle inequality after applying the chosen Gram-frame
isomorphism.

The good-cover size condition gives both
\(\log d_K\le o\) and \(\log\operatorname{rd}K\le o\), with
\(o\le391X^2\).
If the maximum is \(1\), its logarithm is zero.
Otherwise
\[
 \log A\le5o+C_bX^3+C_2(X^3S+X^5)
          \le(2000+C_b+C_2)(X^3S+X^5),
\]
since \(X^2,X^3\le X^5\) and \(5\cdot391<2000\).
For a nonempty family,
\[
 \log n_{\rm gen}\le\log2+\log d_K+\log h
        \le1+391X^2+C_bX^3\le CX^3 .
\]
This is the full numerical passage in
\srcid{FlatConstantAbsLeNineExplicit.short\_gen\_comp}.
Neither the number nor the norm of these generators is obtained
from a Hilbert--Samuel asymptotic.
\end{proof}

\paragraph{Exact scope of the two twists.}
The floor/Minkowski proof is
\srcid{DivisorProofs/ExistsShortIntegralFlatLineTwistingSection.lean}
(\srcid{7a5c6eafdd31c9448b7819342562331dac7b4923}).
The lower-degree identity is
\srcid{DivisorProofs/ExistsTwSecOfConst.lean}
(\srcid{3906b3d7afd777e9d6ec9b67a47e3c394e60176e}).
The multiplicative lattice and norm arguments are
\srcid{ExistsLinearEquivMulTwistingSectionIntegral.lean}
(\srcid{398f8190dbac2615bb6e0612cc9890a80c720b02}) and
\srcid{ExistsSpanningFamilyMulTwistingSection.lean}
(\srcid{8e96b74a7f66eac17ed4b00ae441a570308494a6}).
The exact large-gap exponents are
\srcid{ExistsThetaWeightExponentsOfLargeGap.lean}
(\srcid{af7b06e0d1be0dd845fb0504619cd9b29d6baa52}).
The full assembly is \srcid{ShortLatticeSpanningSets.lean},
followed by \srcid{FlatConstantAbsLeNineExplicit.lean}.

The key-power application is
\srcid{ExistsKeyUnitPowTwistingSectionOnComponent.lean}
(\srcid{d7dd57cc9990ada2ddb8d85b56d821d4a4b86353}).
It requires the actual exponent-one theta rule \(\kappa\), its
functorial family-unit theorem, the componentwise fourth root,
the complex key-norm formula and the finite model-unit/different
comparison.
The first of these is supplied in the frozen source by
\srcid{ThetaLawProofs/KeyFamOne.lean}
and \srcid{ModuliProofs/ExactKeyFormulaProof.lean}, through
\srcid{UniformizationProofs/KeyFormulaOfConnFam.lean}.
The existing component-root/different argument may be reused
only with those actual inputs and the same N3/N4 models.
The cusp-preserving clauses of the twists use
\srcid{MulMemCuspFormsOfBoundedSectionComponent.lean},
including its key-section and analytic hypotheses.
These exact producer joins are tracked separately; the elementary
rounding and weight/numerical deductions above do not by themselves
assert their construction.
